% Modernised reading — fonds Alexandre Grothendieck, shelfmark 98, pages 1-43.
%
% This is an INTERPRETATION, not a transcription. It re-reads the pages in
% today's notation and vocabulary, states the hypotheses the manuscript leaves
% implicit, and moves everything the critical apparatus carried into
% footnotes. It is held to being mathematically correct as it stands; where
% the manuscript is loose or mistaken the true statement appears and the
% footnote says what the page has. In any dispute of substance the
% transcription governs: whoever cites Grothendieck cites batch-01.fr.tex,
% batch-02.fr.tex and batch-03.fr.tex.
%
% Pass: Opus 5.5 (claude-opus-5-5), 2026-10-10 — first pass, unchecked against the pages by a human.
%
% Scope: a single document for the whole folder, read whole (three batches,
% pages 1-43) before any of it was written. Pages 1 and 25 are covers in his
% hand (« Formalisation des Normes / Cf aussi SGA 4 XVIII », « Familles
% limitées / cf. SGA 6 XIII »). Pages 3, 5, 7, 9, 14, 16, 20, 22, 24 are
% leaves of two French typescripts not by him (n° 329, general topology;
% n° 268, abelian varieties), skipped by the transcription; page 18 is one
% such leaf carrying pencil marks, kept there. This reading restates none of
% those typescripts (RIGHTS.md): it says only what the transcription's notes
% say, and less.
%
% The group. The folder is filed in the group « Géométrie et topologie
% combinatoire » (69-102), but its subject and its date ([à partir de
% 1966-1967]) are not the group's; the transcriptions say so and conclude
% nothing from it, and nor does this reading. It does not meet the folders
% of the group already modernised (81, 88, 89: maps, dessins, combinatorial
% hexagons, 1976-1986): no cross-reference to them is made beyond saying so.
%
% Spine. Two halves, matching the inventory's double title « Schémas de
% Hilbert. Normes ». I (2-8) norms: the determinant as a multiplicative
% polynomial law, the canonical section Can_{T/S} : S -> Sym^r_S(T) of a
% finite locally free T -> S of rank r, the norm N^phi(alpha) of a section
% along a symmetric morphism phi, and its transitivity. II (10-13) an
% axiomatics of « morphisms with norms » with three examples (G-sets, finite
% locally free morphisms of schemes, coverings). III (15-23) the cohomology of
% G-sets seen as sheaves on a site, the Čech-to-derived spectral sequences
% (A), (B) of a covering, a low-degree sequence for G_m, and a table of four
% topological fibrations. IV (25-43) limited (bounded) families of coherent
% sheaves: the boundedness of quotients with given Hilbert polynomials
% (25-27), classes of sheaves over extensions of the base field (28-33),
% stability of limited families under ⊗, Hom, Tor, Ext, f^* (34-37),
% base change for Tor and Ext (38-42), the plan of a chapter (43).
%
% Notation fixed for the whole folder. Sym^r_S(T) = T^r_S / S_r (his
% Σ^r(T/S), Symm^n_S); Can_{T/S} the canonical section; N^phi_{T/S}(alpha)
% the norm. 𝔉(X/k) the set of k-classes of coherent sheaves (his gothic F);
% K~ an algebraically closed extension of infinite transcendence degree.
% Sheaf Hom, Tor, Ext are \mathcal{H}om, \mathcal{T}or, \mathcal{E}xt.
% C_G the G-sets with finitely many orbits.
%
% Substantive departures from the pages, each footnoted where it occurs:
%  - p. 2: the action of symmetric endomorphisms on Λ^n M read inside ⊗^n M
%    is literally right only when n! is invertible; the reading states the
%    determinant as a multiplicative polynomial law (Roby), valid over any A.
%  - pp. 4-8: the letters F, G, X, Y change role from sheet to sheet; one
%    convention is fixed (T -> S finite locally free, G the target scheme).
%    The F of the index N^chi_{F/Y/Z} on p. 8 is G''.
%  - p. 12: « conditions (i) à (iv) » for Can read (i ter) to (iv ter).
%  - p. 13: freeness and proper discontinuity of the action of Gamma are
%    supplied (the page is largely illegible there).
%  - p. 15: Φ_A(X) = Hom_G(X, E) read Hom_G(X, A); the topology on C_G is
%    taken to be the one induced from G-Sets (comparison lemma), the page's
%    definition of « hyperfoncteur » being partly illegible.
%  - p. 17: « Hom_Z(Z(G), M) = M^G » is the set of maps G -> M (coinduced
%    module), not the invariants.
%  - p. 21: the exact sequence holds when the row q = 1 of the spectral
%    sequence vanishes; the page does not say so.
%  - p. 23: the local systems are needed in Leray's sequence too; X written
%    for the group Y, B_G for B_X: corrected.
%  - p. 30: « algébrique(ment close) de transcendance … » completed as
%    algebraically closed of infinite transcendence degree.
%  - p. 34: Ḡ = G ⊗_{S'} S''' read ⊗_{S''}; p. 35: Hom(L_i, F) read
%    Hom(L_i, G); p. 36: K^i -> K^{i+2} read K^{i+1}.
%  - p. 36: Tor^{O_Z} read Tor^{O_Y}. p. 37: « Y-plats » read S-flat.
%  - p. 38: the Tor of the Proposition and of the Lemma are relative to Y
%    (the page writes O_{S'}, O_T, O_S); the Lemma needs the F_i flat over S,
%    which its proof (p. 39) uses and its statement omits — supplied.
%  - p. 40: Ext over O_{S'} read over O_{X'}; T -> S read T -> S'.
%  - p. 42: Hom_A read Hom_B; « K^i A-libres » read A-flat; « isom. en tt
%    dimension » read for p <= N.
%
% Announced and never established in the folder: the transitivity of norms
% (p. 8, « Je dis que », no proof); the verification of (i ter)-(iv ter) in
% Example II (« de façon naturelle »); what Φ does to canonical morphisms
% (p. 13, breaks off); « Tout ceci est à généraliser aux hyperfaisceaux
% quelconques » (p. 19); the induction of p. 27 (its step for F'' largely
% illegible); the « Question » of p. 29 (answered by Cor. 5, p. 30, in this
% reading's completion); parts (b) and (c) of the list of p. 33 (illegible);
% the Hom stratification of p. 35 (illegible); the whole of the plan of p. 43
% (Prop. 1-7 by their headings only).
\documentclass[11pt,a4paper]{article}
\input{../preamble/grothendieck}

\folder{98}
\pages{1}{43}
\dating{[à partir de 1966-1967]}
\watermark{Édition de démonstration\\interprétation personnelle de l'œuvre}
\foldertitle{Schémas de Hilbert. Normes : notes manuscrites (s.d.) — lecture modernisée du dossier entier}

\begin{document}

\section*{Résumé}

\begin{resume}
Ces feuillets portent sur deux outils dont on a besoin pour fabriquer des
espaces qui paramètrent des objets géométriques --- ce qu'on appelle des
espaces de modules, dont le schéma de Hilbert est l'exemple fondateur.

Le premier outil est la \emph{norme}. Tout le monde connaît le déterminant
d'une matrice, et la norme d'un nombre complexe, $N(a+bi) = a^{2}+b^{2}$,
qui est le déterminant de la multiplication par $a+bi$ vue comme application
linéaire du plan. Ces normes sont multiplicatives, et transitives : dans une
tour d'extensions, la norme du haut en bas est la composée des normes
intermédiaires. Grothendieck cherche ce qu'il y a de plus général derrière.
Sa réponse : quand un espace $T$ est posé « au-dessus » d'un espace $S$
comme un revêtement à $r$ feuillets --- au-dessus de chaque point de $S$,
$r$ points de $T$, comptés avec multiplicité ---, chaque point de $S$
détermine le paquet non ordonné des $r$ points au-dessus de lui. C'est une
application canonique de $S$ vers la « puissance symétrique » de $T$, l'espace
des paquets de $r$ points. Toute norme, toute trace, se déduit de cette seule
application. Il en dresse les propriétés, puis les érige en axiomes valables
aussi bien pour des ensembles munis d'une action de groupe que pour des
schémas ou des revêtements d'espaces topologiques.

Les ensembles munis d'une action de groupe le mènent à la cohomologie des
groupes, qu'il relit comme la cohomologie d'un espace dont les « ouverts »
sont ces ensembles. Recouvrir un tel espace et comparer la cohomologie du
tout à celle des morceaux donne une suite spectrale, qui contient comme cas
particuliers celles que l'on connaît en topologie --- Leray, Cartan--Leray,
Hochschild--Serre ---, dont il dresse le tableau.

Le second outil est la notion de \emph{famille limitée}. Pour qu'un espace
de modules existe et soit raisonnable, il faut que les objets qu'il
paramètre ne soient pas trop nombreux : qu'on puisse les faire tous tenir
dans une seule famille algébrique, paramétrée par un espace de dimension
finie. Il faut aussi donner un sens précis à « le même faisceau », quand deux
faisceaux vivent sur des corps de base différents --- comme on dirait que
deux équations à coefficients dans des corps différents définissent « la même
courbe ». Il montre ensuite que les opérations usuelles sur les faisceaux
préservent le caractère limité, et pour cela il lui faut savoir quand ces
opérations commutent aux changements de base ; la réponse est la platitude,
et elle se lit sur une suite spectrale. Le dossier se termine sur le plan
d'un chapitre entier consacré aux familles limitées.

Rien de ce qui est ici ne relève de la combinatoire du groupe où le dossier
est rangé : c'est l'atelier de la géométrie algébrique des années 1966--1967,
celui où s'élaborent les séminaires SGA 4 et SGA 6 que les couvertures
citent. Les noms sous lesquels chercher la suite sont ceux de la
section canonique de la puissance symétrique, du morphisme de Hilbert--Chow,
de la norme au sens de Deligne et de Ferrand, des familles bornées de Kleiman.

\keywords{norm map, symmetric power, multiplicative polynomial law, divided powers, finite locally free morphism, Hilbert-Chow morphism, transitivity of norms, Čech-to-derived spectral sequence, Cartan-Leray spectral sequence, Hochschild-Serre spectral sequence, comparison lemma, Brauer group, bounded family, Quot scheme, Chow variety, generic flatness, cohomology and base change, Künneth spectral sequence, relative Ext}
\end{resume}

\section*{Le fil du dossier, et les conventions}

\subsection*{Les stations}

\begin{enumerate}
\item[I.] \emph{Normes} (pages 1 à 8) : le déterminant comme loi
polynôme multiplicative (page 2) ; la section canonique d'un morphisme
fini localement libre dans sa puissance symétrique, et la norme qu'elle
induit (pages 4 à 6) ; la liste des propriétés voulues, et la
transitivité, énoncée sans preuve (pages 6 à 8).
\item[II.] \emph{Théorie axiomatique} (pages 10 à 13) : une catégorie, une
classe de morphismes munis d'un degré et de sections canoniques, quatre
axiomes pour celles-ci ; trois exemples --- $G$-ensembles, schémas,
revêtements.
\item[III.] \emph{Cohomologie des $G$-ensembles et suites spectrales de
recouvrement} (pages 15 à 23) : les $G$-ensembles comme faisceaux sur un
site, les suites spectrales (A) et (B), une suite exacte de bas degré pour
$\mathbf{G}_m$, un tableau de quatre fibrations.
\item[IV.] \emph{Familles limitées} (pages 25 à 43) : bornitude des
quotients d'un faisceau fixe (pages 25 à 27) ; classes de faisceaux sur les
extensions du corps de base (pages 28 à 33) ; stabilité des familles
limitées (pages 34 à 37) ; changement de base pour les $\mathcal{T}or$ et les
$\mathcal{E}xt$ (pages 38 à 42) ; plan d'un chapitre (page 43).
\end{enumerate}

Les deux moitiés répondent aux deux mots du titre de l'inventaire,
« Normes » et « Schémas de Hilbert ». Le dossier ne dit pas ce qui les
relie. Une liaison naturelle, qui est de cette lecture et non des pages, est
le morphisme de Hilbert--Chow : sur le schéma de Hilbert des sous-schémas
finis de degré $r$ de $X/S$, la famille universelle est finie localement
libre de rang $r$, et sa section canonique (station I) donne un morphisme
vers la puissance symétrique $\mathrm{Sym}^r_S(X)$.

Les feuillets de la station III sont reliés aux premières par l'exemple I
de la page 11 : la catégorie des $G$-ensembles y sert de modèle à
l'axiomatique, et c'est sur elle (pages 15 à 19) que s'écrit la
cohomologie.

\subsection*{Conventions}

Tous les anneaux sont commutatifs et unitaires. Un morphisme $p : T \to S$
est \emph{fini localement libre de rang $r$} si $p_{*}\mathcal{O}_T$ est un
$\mathcal{O}_S$-module localement libre de rang $r$. On note
$T^{r}_{S} = T \times_S \cdots \times_S T$ ($r$ facteurs), et
\[
  \mathrm{Sym}^{r}_{S}(T) = T^{r}_{S}/\mathfrak{S}_r
\]
sa puissance symétrique, qui existe dès que $T$ est affine sur $S$ : si
$S = \operatorname{Spec} A$ et $T = \operatorname{Spec} B$, c'est
$\operatorname{Spec} \mathrm{TS}^{r}_{A}(B)$, où
$\mathrm{TS}^{r}_{A}(B) = (B^{\otimes r})^{\mathfrak{S}_r}$ est l'algèbre des
tenseurs symétriques. La page écrit tour à tour $\Sigma^{r}(T/S)$ et
$\mathrm{Symm}^{r}_{S}(T)$ ; on écrit $\mathrm{Sym}^{r}_{S}(T)$ partout.
La section canonique est $\mathrm{Can}_{T/S} : S \to \mathrm{Sym}^{r}_{S}(T)$,
et la norme d'une section $\alpha$ le long d'un morphisme symétrique
$\varphi$ est $N^{\varphi}_{T/S}(\alpha)$.\footnote{Les lettres changent de
rôle d'une feuille à l'autre des pages 4 à 8 : le revêtement s'appelle tour à
tour $T \to S$, $X \to S$, $X \to Y$, et l'indice de la norme porte trois
lettres, $N^{\alpha}_{F/T/S}$, la première nommant le schéma but. On fixe ici
une seule convention, et l'on dit en note, à chaque fois, ce que la page
écrit.}

Pour la seconde moitié, $S$ est un schéma noethérien, $X$ un $S$-schéma de
type fini ; les faisceaux sont cohérents. Les foncteurs faisceautiques sont
notés $\mathcal{H}om$, $\mathcal{T}or$, $\mathcal{E}xt$ (il les souligne).
$\mathfrak{F}(X/k)$ est l'ensemble des classes de faisceaux cohérents de la
page 29, et $\widetilde{K}$ une extension algébriquement close de $k$ de
degré de transcendance infini. Les « préschémas » des pages sont nos schémas.

Le mot « limité » est le sien, et celui de SGA 6 ; on dit aujourd'hui
« borné » (\emph{bounded}). On garde « limité », en le rappelant.

\pagerange{1}{2}
\section*{I. Normes (pages 1 à 8)}

La couverture porte, de sa main, « Formalisation des Normes » et
« Cf aussi SGA 4 XVIII ».\footnote{C'est dans SGA 4, chez Deligne, que la
trace d'un morphisme fini localement libre est construite au moyen de la
puissance symétrique (exposé XVII, n° 6.3, à notre connaissance) ;
l'exposé XVIII, que cite la couverture, est celui de la dualité globale,
qui s'en sert. On ne prétend pas que ces feuillets précèdent ou suivent ces
exposés : la datation de l'inventaire, « à partir de 1966-1967 », ne le
tranche pas.}

\subsection*{Le déterminant comme loi multiplicative (page 2)}

Soient $A$ un anneau et $M$ un $A$-module libre de rang $n$. On a
\[
  L_A\bigl(T^{n}_{A}(M)\bigr) \simeq T^{n}_{A}\bigl(L_A(M)\bigr) \supset
  \mathrm{TS}^{n}_{A}\bigl(L_A(M)\bigr),
\]
où $L_A$ désigne les endomorphismes et $T^{n}$ la puissance tensorielle :
les tenseurs symétriques de $L_A(M)$ sont exactement les endomorphismes de
$\otimes^{n} M$ qui commutent à l'action de $\mathfrak{S}_n$. Pour
$u \in L_A(M)$, $u^{\otimes n}$ en est un, et il opère sur la puissance
extérieure maximale $\Lambda^{n} M$, libre de rang 1, par la multiplication
par $\det u$.

L'énoncé juste, sur un anneau quelconque, est le suivant : l'application
$u \mapsto \det u$ est une loi polynôme homogène de degré $n$ et
multiplicative de $L_A(M)$ dans $A$ ; elle correspond donc à un
homomorphisme d'algèbres
\[
  \mathrm{TS}^{n}_{A}\bigl(L_A(M)\bigr) \simeq \Gamma^{n}_{A}\bigl(L_A(M)\bigr)
  \longrightarrow A,
\]
où $\Gamma^{n}$ est la puissance divisée, isomorphe aux tenseurs
symétriques pour un module libre.\footnote{La page fait opérer les
endomorphismes symétriques sur « la partie stable par symétrie » de
$\otimes^{n} M$ qu'elle identifie à $\Lambda^{n} M$. C'est exact quand $n!$
est inversible dans $A$ : $\Lambda^{n} M$ est alors le facteur direct des
tenseurs anti-invariants. En général --- en caractéristique 2, les tenseurs
anti-invariants sont les tenseurs symétriques ---, il faut passer par les
lois polynômes de Roby (1963), dont la page ne parle pas ; l'énoncé ci-dessus
est le nôtre. La correspondance entre lois polynômes multiplicatives
homogènes de degré $n$ et homomorphismes d'algèbres
$\Gamma^{n}_{A}(B) \to A$ est de Roby.}

Si $B$ est une $A$-algèbre finie localement libre de rang $n$, en composant
avec $B \to L_A(B)$, $b \mapsto (x \mapsto bx)$, on obtient la norme
$N_{B/A}$ comme loi multiplicative de degré $n$, c'est-à-dire un
homomorphisme d'algèbres $\mathrm{TS}^{n}_{A}(B) \to A$ : géométriquement,
une section
\[
  \mathrm{Can}_{T/S} : S \longrightarrow \mathrm{Sym}^{n}_{S}(T),
  \qquad T = \operatorname{Spec} B,\ S = \operatorname{Spec} A .
\]
Elle se recolle, et définit $\mathrm{Can}_{T/S}$ pour tout $T \to S$ fini
localement libre de rang $n$. C'est le passage, au bas de la page, de
$\otimes^{n} F$ à $X^{n}_{S} \to \mathrm{Symm}^{n}_{S}(Y)$. La propriété
universelle du quotient y est esquissée : un morphisme $Y^{n}_{S} \to Z$
invariant par $\mathfrak{S}_n$ se factorise de façon unique par
$\mathrm{Sym}^{n}_{S}(Y)$.\footnote{Le diagramme de la page 2 met en regard
$X^{n}_{S} \to Z$, $Y^{n}_{S} \to \mathrm{Symm}^{n}_{S}(Y) \to Z$, et des
flèches verticales $Y \to X$, $Y^{n}_{S} \to X^{n}_{S}$ dont le rôle n'est
pas dit ; on le lit comme la propriété universelle du quotient, ce qui est
une interprétation. Il écrit $\mathrm{Cart}^{n}_{S}(X)$ et
$\overset{n}{\boxtimes}_S X$ pour le produit fibré itéré.}

\pagerange{4}{6}
\subsection*{La norme le long d'un morphisme symétrique (pages 4 à 6)}

Soit $p : T \to S$ fini localement libre de rang $r$. Soient $G$ et $G'$
deux $S$-schémas et $\varphi : G^{r}_{S} \to G'$ un $S$-morphisme
\emph{symétrique}, c'est-à-dire invariant par $\mathfrak{S}_r$. Une section
de $G_T = G \times_S T$ au-dessus de $T$ est la même chose qu'un
$S$-morphisme $\alpha : T \to G$. Le morphisme
$\varphi \circ \alpha^{r} : T^{r}_{S} \to G'$ est symétrique, donc se
factorise par $\mathrm{Sym}^{r}_{S}(T)$ ; on pose
\[
  N^{\varphi}_{T/S}(\alpha) : S \xrightarrow{\ \mathrm{Can}_{T/S}\ }
  \mathrm{Sym}^{r}_{S}(T) \xrightarrow{\ \overline{\varphi \circ \alpha^{r}}\ } G' .
\]
On a ainsi une application
\[
  N^{\varphi}_{T/S} : \Gamma(G_T/T) \longrightarrow \Gamma(G'/S).
\]
Pour $G = G' = \mathbf{A}^{1}_{S}$ et $\varphi$ le produit des $r$
coordonnées, c'est la norme usuelle $\Gamma(T, \mathcal{O}_T) \to
\Gamma(S, \mathcal{O}_S)$.\footnote{La page 4 écrit
$N^{\alpha}_{F/T/S} : \Gamma(G/T) \to \Gamma(F/S)$, où $\alpha$ est le
morphisme symétrique, $G$ le $T$-schéma, $F$ le $S$-schéma but ; une
première forme, $N^{F}_{T/S} : \Gamma(G_T/T) \to \Gamma(G/S)$, la précède.
La page ne dit nulle part que la puissance symétrique de $G$ existe, et notre
définition n'en a pas besoin : on ne forme que celle de $T$, qui est affine
sur $S$.}

La page 6 dresse la liste des propriétés à établir pour $\mathrm{Can}$ :
\begin{enumerate}
\item[a)] les cas $r = 0$ et $r = 1$ : $\mathrm{Sym}^{0}_{S}(T) = S$, et
pour $r = 1$, $T = S$ et $\mathrm{Can}_{S/S} = \mathrm{id}_S$ ;
\item[b)] la fonctorialité : pour $S' \to S$ et $T' = T \times_S S'$, la
section $\mathrm{Can}_{T'/S'}$ est l'image inverse de $\mathrm{Can}_{T/S}$, et
la norme commute aux morphismes $G \to G_1$, $G' \to G'_1$ qui rendent
commutatifs les carrés évidents ;
\item[c)] l'additivité : si $T = T_1 \amalg T_2$, de rangs $r_1$, $r_2$,
alors $\mathrm{Can}_{T/S}$ est le composé de
$(\mathrm{Can}_{T_1/S}, \mathrm{Can}_{T_2/S})$ et du morphisme canonique
$\mathrm{Sym}^{r_1}_{S}(T_1) \times_S \mathrm{Sym}^{r_2}_{S}(T_2) \to
\mathrm{Sym}^{r_1+r_2}_{S}(T)$ ;
\item[d)] la transitivité, objet de la page 8.
\end{enumerate}
Pour les normes de fonctions, c) dit que la norme d'une fonction sur
$T_1 \amalg T_2$ est le produit des normes de ses restrictions.\footnote{La
page intitule c) « Additivité en $M$ » : le module $M$ de la page 2,
c'est-à-dire $p_{*}\mathcal{O}_T = M_1 \oplus M_2$. Dans ses diagrammes, le
membre de gauche est $X^{r_1+r_2}_{S}$ et la flèche verticale
$X^{r_1}_{S} \times_S X^{r_2}_{S} \to X^{r_1+r_2}_{S}$ est marquée
« canon. » et d'un $\wr$. Le symbole lu $\mathrm{Can}$ est écrit ici presque
« Cm » ; la page 10 le fixe.}

\pagerange{6}{8}
\subsection*{La transitivité (pages 6 à 8)}

Soient $X \xrightarrow{u} Y \xrightarrow{v} Z$ finis localement libres de
rangs $r$ et $s$ ; le composé est de rang $rs$. Il y a un morphisme
canonique
\[
  \mathrm{can} : (X^{r}_{Y})^{s}_{Z} \longrightarrow X^{rs}_{Z},
\]
qui envoie $s$ paquets de $r$ points, chacun au-dessus d'un point de $Y$, sur
leur juxtaposition ; il passe aux quotients par
$\mathfrak{S}_r \wr \mathfrak{S}_s \subset \mathfrak{S}_{rs}$.

Soient $G$ un $Z$-schéma, et donnons-nous :
\begin{enumerate}
\item[--] un $Z$-morphisme symétrique $\varphi : G^{rs}_{Z} \to G'$ ;
\item[--] un $Y$-schéma $G''$ et un $Y$-morphisme symétrique
$\psi : (G_Y)^{r}_{Y} \to G''$ ;
\item[--] un $Z$-morphisme symétrique $\chi : G''^{\,s}_{Z} \to G'$, où
$G''$ est vu comme $Z$-schéma,
\end{enumerate}
tels que le carré
\[
  \chi \circ \psi^{s} = \varphi \circ \mathrm{can} :
  \bigl((G_Y)^{r}_{Y}\bigr)^{s}_{Z} \longrightarrow G'
\]
commute. Soit $\alpha : X \to G$ un $Z$-morphisme, et
$\alpha' : X \to G_Y$ le $Y$-morphisme qui lui correspond
($\alpha = c \circ \alpha'$, $c : G_Y \to G$ la projection). La page affirme,
encadré :
\[
  \boxed{\;N^{\varphi}_{X/Z}(\alpha) =
  N^{\chi}_{Y/Z}\bigl(N^{\psi}_{X/Y}(\alpha')\bigr)\;}
\]
où $N^{\psi}_{X/Y}(\alpha') : Y \to G''$ est vu comme une section de
$G''$ au-dessus de $Y$, c'est-à-dire un $Z$-morphisme $Y \to G''$.

C'est la transitivité des normes : pour $G = G' = G'' = \mathbf{A}^{1}$ et
les produits, elle redonne $N_{C/A} = N_{B/A} \circ N_{C/B}$ pour
$A \to B \to C$ finis localement libres. Elle découle de la transitivité de
la section canonique, qui est l'axiome (ii ter) de la page 11 ci-dessous ;
la page la donne par « Je dis que », sans preuve.\footnote{Pages 4 et 8.
Sur la page 4, elle est écrite
$N^{\gamma}_{(G \otimes_T F)/U/S}(\varphi) = N^{\alpha}_{F/T/S}(N^{\beta}_{G/U/T}(\varphi))$,
avec $U \to T \to S$ de rangs $s$, $r$, et $\varphi$ pour la section. Sur la
page 8 l'indice du membre de droite est $F/Y/Z$, écrit par-dessus $F/X/Z$ :
le schéma but y est notre $G''$. La page 6 encadre les quatre données sous
la forme $X^{r}_{Y} \to \mathfrak{X}_r$, $Y^{s}_{Z} \to \mathfrak{Y}_s$,
$X^{r}_{Z} \to \mathfrak{Z}_r$, $\mathfrak{X}^{s}_{Z} \to \mathfrak{Z}$, avec
« Commutativité » : un exposant lu $r$ y est peut-être $rs$. Plusieurs mots
de la page 8 (« suppose », « rendant », « appliquant », « trouve ») sont
d'une lecture incertaine ; le sens des formules ne l'est pas. Une preuve de
la transitivité de la norme dans cette généralité, pour les algèbres finies
localement libres, se trouve, à notre connaissance, dans le « foncteur
norme » de D.~Ferrand (1998) ; la page ne la donne pas.}

\pagerange{10}{11}
\section*{II. Une théorie axiomatique des normes (pages 10 à 13)}

\subsection*{Les axiomes (pages 10 et 11)}

Soit $\mathcal{C}$ une catégorie à sommes finies, et $\mathcal{N}$ une
classe de morphismes de $\mathcal{C}$ telle que :
\begin{enumerate}
\item[(i)] les identités sont dans $\mathcal{N}$ ;
\item[(ii)] $\mathcal{N}$ est stable par composition ;
\item[(iii)] si $u : X \to Y$ est dans $\mathcal{N}$ et $Y' \to Y$
quelconque, $X' = X \times_Y Y'$ existe et $u' : X' \to Y'$ est dans
$\mathcal{N}$ ;
\item[(iv)] si $u_1$, $u_2$ sont dans $\mathcal{N}$, de même but $Y$,
$(u_1, u_2) : X_1 \amalg X_2 \to Y$ y est aussi ;
\item[(v)] si $u : X \to Y$ est dans $\mathcal{N}$, pour tout entier $n$ et
tout sous-groupe $H \subset \mathfrak{S}_n$, le quotient
$(X^{n}_{Y})/H$ existe (et sa projection sur $Y$ est dans
$\mathcal{N}$).
\end{enumerate}
On se donne un degré $\deg : \mathcal{N} \to \mathbf{N}$ avec
\begin{enumerate}
\item[(i bis)] $\deg(\mathrm{id}_X) = 1$ ;
\item[(ii bis)] $\deg(vu) = \deg v \cdot \deg u$ ;
\item[(iii bis)] le degré est invariant par changement de base ;
\item[(iv bis)] $\deg(u_1, u_2) = \deg u_1 + \deg u_2$ ;
\item[(v bis)] si $\deg u = r$ et $E = \{1, \ldots, r\}$, le degré de
$(X^{n}_{Y})/H \to Y$ est $\mathrm{card}(E^{n}/H)$.
\end{enumerate}
Enfin, pour $u : X \to Y$ dans $\mathcal{N}$ de degré $r$, on se donne
$\mathrm{Can}_{X/Y} : Y \to \mathrm{Sym}^{r}_{Y}(X)$, où
$\mathrm{Sym}^{r}_{Y}(X) = (X^{r}_{Y})/\mathfrak{S}_r$, avec :
\begin{enumerate}
\item[(i ter)] $\mathrm{Can}_{X/X} = \mathrm{id}_X$ ;
\item[(ii ter)] pour $X \xrightarrow{u} Y \xrightarrow{v} Z$ de degrés $r$,
$s$, le diagramme ci-dessous commute ;
\item[(iii ter)] $\mathrm{Can}$ commute au changement de base $Y' \to Y$ ;
\item[(iv ter)] pour $X = X_1 \amalg X_2$, on a
$\mathrm{Sym}^{r}_{Y}(X) = \coprod_{i+j=r} \mathrm{Sym}^{i}_{Y}(X_1)
\times_Y \mathrm{Sym}^{j}_{Y}(X_2)$, et $\mathrm{Can}_{X/Y}$ est le composé
de $(\mathrm{Can}_{X_1/Y}, \mathrm{Can}_{X_2/Y})$ et de l'inclusion du
facteur d'indice $(r_1, r_2)$.
\end{enumerate}
Le diagramme de (ii ter) est
\[
  Z \xrightarrow{\ \mathrm{Can}_{Y/Z}\ } \mathrm{Sym}^{s}_{Z}(Y)
  \xrightarrow{\ \mathrm{Sym}^{s}(\mathrm{Can}_{X/Y})\ }
  \mathrm{Sym}^{s}_{Z}\bigl(\mathrm{Sym}^{r}_{Y}(X)\bigr)
\]
\[
  \xrightarrow{\ \mathrm{can}\ }
  \mathrm{Sym}^{s}_{Z}\bigl(\mathrm{Sym}^{r}_{Z}(X)\bigr)
  \xrightarrow{\ \mathrm{can}\ } \mathrm{Sym}^{rs}_{Z}(X),
\]
dont le composé doit être $\mathrm{Can}_{X/Z}$.\footnote{Ce sont les
axiomes de la page, aux mots près. La fin de (iv) est très corrigée et
illisible ; celle de (v) aussi, entourée au crayon rouge d'un « ? » : on
complète par ce que (v bis) suppose, que le quotient est dans $\mathcal{N}$,
et c'est notre ajout. La page écrit $\Sigma^{r}(X/Y)$. Dans (v bis) elle
emploie $r$ et $n$ à la fois ; le groupe y est appelé $G$, qu'on renomme $H$
pour ne pas le confondre avec le groupe de l'exemple I. Le terme du milieu
de la colonne de droite du diagramme de la page 11 est écrit sur un premier
terme, peut-être $\Sigma^{s+r}$. Dans (iv ter) la page écrit $\times$ pour
$\times_Y$.}

\pagerange{11}{13}
\subsection*{Trois exemples (pages 11 à 13)}

\emph{Exemple I : les $G$-ensembles.} $G$ est un groupe, $\mathcal{C}$ la
catégorie des ensembles où $G$ opère ; $u : X \to Y$ est de degré $r$ si
toutes ses fibres ont $r$ éléments, et $\mathcal{N}$ est formé des
morphismes qui ont un degré. Produits fibrés et quotients par un groupe
d'automorphismes existent dans $\mathcal{C}$, et les axiomes (i) à (v),
(i bis) à (v bis) sont immédiats. L'ensemble $\mathrm{Sym}^{r}_{Y}(X)$ est
la réunion disjointe des $\mathrm{Sym}^{r}(u^{-1}(y))$, $y \in Y$ ; comme
$u^{-1}(y)$ a exactement $r$ éléments, $\mathrm{Sym}^{r}(u^{-1}(y))$ a un
élément privilégié, la classe de la suite de ses éléments rangés dans un
ordre quelconque. D'où $\mathrm{Can}_{X/Y}$, évidemment compatible à
l'action de $G$, et (i ter) à (iv ter) se vérifient aussitôt.\footnote{Le
bas du feuillet droit de la page 11 est en partie hors de l'image, et
plusieurs mots de l'argument sont illisibles ; la conclusion se lit en tête
de la page 12.}

\emph{Exemple II : les schémas.} $\mathcal{C}$ est la catégorie des
schémas, $\mathcal{N}$ la classe des morphismes finis localement libres de
rang constant, le degré étant le rang. Les axiomes se vérifient « très
facilement », dit la page, et c'est vrai pour la raison suivante : si
$Y = \operatorname{Spec} A$ et $X = \operatorname{Spec} B$ avec $B$ libre
de base $(e_1, \ldots, e_r)$, alors $B^{\otimes n}$ est un module de
permutation, de base $E^{n}$, sur lequel $H$ opère en permutant la base ;
ses invariants sont libres, de base les sommes d'orbites, de rang
$\mathrm{card}(E^{n}/H)$, et ils commutent à tout changement de base. On
définit $\mathrm{Can}_{X/Y}$ par la norme de la page 2.\footnote{La page dit
« fini et plat … loc. libre de rang $r$ en tout point » : on prend « fini
localement libre », qui est ce que dit la fin de la phrase et qui coïncide
avec « fini plat de présentation finie ». Elle affirme que $\mathrm{Can}$
vérifie « les conditions (i) à (iv) » : ce sont (i ter) à (iv ter). La
vérification n'est pas donnée ; l'argument par les modules de permutation
est le nôtre. La note marginale reliée à l'exemple (« Expliquer … relation
… $\mathcal{C}$ … en général … géométriques spéciaux ») est trop lacunaire
pour être interprétée.}

\emph{Exemples III : variantes.} En géométrie analytique, variante évidente.
En topologie, les revêtements finis (non ramifiés) donnent une variante
« très voisine de l'exemple I », que la page 13 précise ainsi. Soit
$X \in \mathcal{C}$ muni d'un groupe discret $\Gamma$ d'automorphismes qui
opère librement et proprement, de sorte que les quotients $X/H$ existent pour
tout sous-groupe $H$, et que pour $H' \subset H$ distingué d'indice fini,
$X/H' \to X/H = (X/H')/(H/H')$ soit un revêtement galoisien de groupe
$H/H'$, donc dans $\mathcal{N}$. Soit $\mathcal{C}'$ la catégorie des
$\Gamma$-ensembles $E$ tels que $E/\Gamma$ soit fini, et
\[
  \Phi(E) = X \times^{\Gamma} E
\]
le produit contracté. Alors $\Phi : \mathcal{C}' \to \mathcal{C}$ commute aux
sommes finies et transforme les morphismes de degré $r$ en morphismes de
degré $r$ ; la page ajoute qu'il transforme les morphismes canoniques en
morphismes canoniques, et s'interrompt.\footnote{Page 13, presque toute la
prose est illisible. Les hypothèses « librement et proprement » sont
fournies par nous : c'est ce qu'il faut pour que $X/H' \to X/H$ soit un
revêtement, et c'est ce que la page désigne par « non ramifié », souligné.
La lettre du groupe y passe de $\Gamma$ à $G$ et $H$. La note marginale
« ; dualement », cerclée, est d'une lecture douteuse ; elle revient aux
pages 15 et 19. C'est, dans le langage de SGA 1, le foncteur qui réalise la
correspondance de Galois entre ensembles à opérateurs et revêtements ; la
page ne le nomme pas.}

\pagerange{15}{15}
\section*{III. Cohomologie des $G$-ensembles et suites spectrales de recouvrement (pages 15 à 23)}

\subsection*{Les $G$-ensembles comme faisceaux (page 15)}

Soit $G$ un groupe, et $\mathcal{C}_G$ la catégorie des ensembles où $G$
opère à gauche et qui n'ont qu'un nombre fini d'orbites. Munissons
$\mathcal{C}_G$ de la topologie induite par la topologie canonique du topos
$B_G$ des $G$-ensembles : les familles couvrantes sont les familles
conjointement surjectives. Comme $\mathcal{C}_G$ est une sous-catégorie
pleine génératrice de $B_G$, le lemme de comparaison donne une équivalence
entre la catégorie des faisceaux sur $\mathcal{C}_G$ --- ses
« hyper-foncteurs » --- et celle de tous les $G$-ensembles. Au
$G$-ensemble $A$ correspond le faisceau
\[
  \Phi_A(X) = \operatorname{Hom}_G(X, A),
\]
et au faisceau $\Phi$ le $G$-ensemble $A = \Phi(G_s)$, où $G_s$ est $G$
opérant sur lui-même par translations à gauche ; l'action de $G$ sur
$\Phi(G_s)$ vient des translations à droite, qui sont les automorphismes de
$G_s$, et elle est à gauche parce que $\Phi$ est
contravariant.\footnote{La page définit les hyper-foncteurs comme des
foncteurs sur $\mathcal{C}_G$ (les mots qui disent ce qu'ils font des
sommes sont illisibles) pour lesquels les épimorphismes de $\mathcal{C}_G$
sont des morphismes de descente. On la lit comme la condition de faisceau
pour la topologie induite, et l'appel au lemme de comparaison (SGA 4, III,
4.1) est le nôtre : il évite une difficulté que la page ne voit pas, à
savoir que $\mathcal{C}_G$ n'est pas stable par produits fibrés
($G/H \times G/K$ a autant d'orbites que $H \backslash G / K$). « À droite »
est biffé et remplacé par « à gauche », deux fois. La formule porte
$\operatorname{Hom}_G(X, E)$, avec $E$ pour $A$. Le mot « hyper-foncteur »
est le sien ; on dit aujourd'hui faisceau sur le site. Un dernier alinéa,
barré de trois traits, introduisait la sous-catégorie des $X$ où un
sous-groupe $F$ opère trivialement.}

\pagerange{17}{18}
\subsection*{La suite spectrale (A) (page 17)}

Soit $A$ un $G$-module, et pour $S \in \mathcal{C}_G$ posons
\[
  \mathcal{H}^{q}(A)(S) = R^{q}\operatorname{Hom}_G(S, -)\,(A)
  = \operatorname{Ext}^{q}_{\mathbf{Z}[G]}\bigl(\mathbf{Z}[S], A\bigr),
\]
de sorte que, par le lemme de Shapiro,
\[
  \mathcal{H}^{q}(A)(F \backslash G) = H^{q}(F, A)
\]
pour tout sous-groupe $F$.\footnote{La page écrit l'orbite $F \backslash G$ ;
pour une action à gauche c'est $G/F$, comme le dit d'ailleurs la formule
biffée qui précède, et comme l'écrit la page 19. On garde la forme de la
page.} C'est un préfaisceau sur $\mathcal{C}_G$ qui transforme sommes
finies en produits. Pour $T \in \mathcal{C}_G$ non vide, notons
$H^{p}(T/e, P)$ la cohomologie de Čech d'un tel préfaisceau $P$ relativement
au recouvrement $T \to e$ de l'objet final $e = G/G$, c'est-à-dire la
cohomologie du complexe
$P(T) \to P(T \times T) \to P(T \times T \times T) \to \cdots$. On a la suite
spectrale
\[
  \text{(A)}\qquad
  E_2^{p,q} = H^{p}\bigl(T/e, \mathcal{H}^{q}(A)\bigr) \Longrightarrow
  H^{p+q}(G, A) \qquad (T \neq \emptyset).
\]
C'est la suite spectrale de Čech d'un recouvrement --- la page dit en marge
« la suite spectrale de Leray d'un \emph{recouvrement} d'un objet » ---, et
elle résulte des deux faits que la page énonce :
\begin{enumerate}
\item[(i)] $T \to e$ est un épimorphisme, donc un morphisme de descente :
pour un faisceau $\Phi$, $\Phi(e) = H^{0}(T/e, \Phi)$ ;
\item[(ii)] si $\Phi$ est injectif, $H^{p}(T/e, \Phi) = 0$ pour $p > 0$.
\end{enumerate}
Pour (ii) il suffit de traiter un module coïnduit
$A = \operatorname{Hom}_{\mathbf{Z}}(\mathbf{Z}[G], M)$, puisque tout
injectif en est facteur direct : alors $\Phi_A(S)$ est l'ensemble des
applications de $S$ dans $M$, et le complexe de Čech est celui des
applications $T^{p+1} \to M$, acyclique en degrés $> 0$ dès que $T$ n'est pas
vide.\footnote{La page écrit
$A = \operatorname{Hom}_{\mathbf{Z}}(\mathbf{Z}(G), M) = M^{G}$ : $M^{G}$
désigne ici l'ensemble des applications de $G$ dans $M$, non les invariants.
La parenthèse de (ii) ne se ferme pas. Le mot « facteur direct » est notre
complément. Les deux $\operatorname{Hom}$ de la définition de
$\mathcal{H}^{q}(A)(S)$ sont biffés sur la page ; ce qui reste est
$R^{q}_{A}(\ldots(S, A))$, que l'on lit comme ci-dessus.}

Le verso de cette feuille, la page 18, est un feuillet d'un tapuscrit qui
n'est pas de lui, sur les variétés abéliennes ; il porte quelques marques au
crayon, dont « Prouver ! », que rien sur la feuille n'attribue à une autre
main, mais qu'un relecteur devra confirmer. Elles ne touchent pas à ces
notes.

\pagerange{19}{19}
\subsection*{La suite spectrale (B) de Leray--Čech (page 19)}

Plus généralement, pour un faisceau abélien $\Phi$ sur $\mathcal{C}_G$ et
$S \in \mathcal{C}_G$, posons
\[
  \mathcal{H}^{q}(\Phi)(S) = H^{q}(-/S, \Phi),
\]
valeur en $\Phi$ du $q$-ième dérivé du foncteur $\Phi \mapsto \Phi(S)$,
c'est-à-dire la cohomologie de la restriction de $\Phi$ au site localisé
en $S$. On montre :
\begin{enumerate}
\item[(i)] $\mathcal{H}^{q}(\Phi)$ est un préfaisceau qui transforme sommes
finies en produits ;
\item[(ii)] $\mathcal{H}^{q}(\Phi_A)(G/F) = H^{q}(F, A)$.
\end{enumerate}
Si $T \to S$ est surjectif dans $\mathcal{C}_G$, on a la suite spectrale
\[
  \text{(B)}\qquad
  E_2^{p,q} = H^{p}\bigl(T/S, \mathcal{H}^{q}(\Phi)\bigr) \Longrightarrow
  H^{p+q}(-/S, \Phi),
\]
que la page appelle suite spectrale de Leray--Čech ; (A) en est le cas
$S = e$. Pour $F' \subset F \subset G$, $T = G/F'$ et $S = G/F$, elle donne,
pour tout $G$-module $A$,
\[
  E_2^{p,q} = H^{p}\bigl(F/F', \mathcal{H}^{q}(A)\bigr) \Longrightarrow
  H^{p+q}(F, A),
\]
la cohomologie de Čech du recouvrement $F/F' \to \mathrm{pt}$ des
$F$-ensembles ; elle n'utilise de $A$ que sa structure de
$F$-module.\footnote{La page écrit $H^{p}(F \bmod F' ; \mathcal{H}^{q}(A))$,
avec un exposant illisible après le premier $F$, et ses deux $A$ sont
surchargés. La phrase qui suit est en partie illisible ; on en garde ce qui
se lit. La note marginale de la formule (B), encadrée, est « suite spectrale
de Leray-Čech », le « Čech » d'une lecture incertaine.} Lorsque $F'$ est
distingué dans $F$, le complexe de Čech de ce recouvrement galoisien est le
complexe standard du groupe $F/F'$ à coefficients dans
$\mathcal{H}^{q}(A)(G/F') = H^{q}(F', A)$, et l'on retrouve la suite de
Hochschild--Serre $H^{p}(F/F', H^{q}(F', A)) \Rightarrow
H^{p+q}(F, A)$.\footnote{Cette identification est de nous ; la page ne
nomme pas Hochschild--Serre ici, mais elle le fait page 23.}

La page conclut : « Tout ceci est à généraliser aux hyperfaisceaux
quelconques ». C'est ce qu'est devenue la suite spectrale de Čech pour un
recouvrement d'un objet d'un topos quelconque, et sa généralisation par les
hyperrecouvrements de Verdier (SGA 4, V).\footnote{Le rapprochement avec
SGA 4, V est le nôtre ; on ne sait pas si « hyperfaisceaux » est ici un
autre nom des hyper-foncteurs de la page 15 ou désigne autre chose.}

\pagerange{21}{21}
\subsection*{Une suite exacte pour $\mathbf{G}_m$ (page 21)}

Appliquée au faisceau $\mathbf{G}_m$ et à un recouvrement $T \to S$, la
suite spectrale de Čech
$H^{p}(T/S, \mathcal{H}^{q}(\mathbf{G}_m)) \Rightarrow H^{p+q}(-/S,
\mathbf{G}_m)$ donne, \emph{lorsque sa ligne $q = 1$ est nulle} --- par
exemple si $\mathcal{H}^{1}(\mathbf{G}_m) = \mathrm{Pic}$ s'annule sur les
puissances fibrées de $T$ ---, la suite exacte
\[
  0 \to H^{2}(T/S, \mathbf{G}_m) \to H^{2}(-/S, \mathbf{G}_m) \to
  H^{0}\bigl(T/S, \mathcal{H}^{2}(\mathbf{G}_m)\bigr) \xrightarrow{\ d_3\ }
  H^{3}(T/S, \mathbf{G}_m) \to H^{3}(-/S, \mathbf{G}_m)'
\]
\[
  \to H^{1}\bigl(T/S, \mathcal{H}^{2}(\mathbf{G}_m)\bigr) \xrightarrow{\ d_3\ }
  H^{4}(T/S, \mathbf{G}_m) \to E_4^{4,0} \to 0,
\]
où $H^{3}(-/S, \mathbf{G}_m)'$ est le noyau de $H^{3}(-/S, \mathbf{G}_m) \to
E_2^{0,3}$, et où $H^{p}(T/S, \mathbf{G}_m)$ est la cohomologie de Čech à
valeurs dans $\mathbf{G}_m$. C'est l'outil qui décrit le groupe de Brauer
par des cocycles de Čech.\footnote{La page commence la suite par « $\cdot
\to$ » et ne dit pas l'hypothèse sur la ligne $q = 1$ ; sans elle la suite
n'est pas exacte telle qu'écrite. Le zéro initial, le sens de
$H^{3}(-/S, \mathbf{G}_m)'$ et l'hypothèse sont les nôtres. Au pied de la
page figure un carré de corps $k \subset k', k''$, $k'k''$, deux côtés
marqués « rad » (radiciel), deux d'une marque lue « sép », et, isolé,
$H^{2}(\ /k, \mathbf{G}_m)$, c'est-à-dire le groupe de Brauer de $k$. On ne
reconstruit pas l'argument que ce carré annonce.}

\pagerange{23}{23}
\subsection*{Quatre fibrations (page 23)}

Le modèle topologique de ces suites spectrales est dressé en tableau.
\begin{enumerate}
\item[--] \emph{Leray.} Pour une fibration $F \to X \to Y$,
$E_2^{p,q} = H^{p}(Y, \mathcal{H}^{q}(F)) \Rightarrow H^{p+q}(X)$, où
$\mathcal{H}^{q}(F)$ est le système local des $H^{q}$ des fibres.
\item[--] \emph{Cartan--Leray.} Pour un fibré principal $X \to Y$ de groupe
$F$, classifié par $Y \to B_F$, de fibre homotopique $X$ :
$H^{p}(B_F, \{H^{q}(X)\}) \Rightarrow H^{p+q}(Y)$, le système local étant
défini par l'action de $\pi_0(F) = \pi_1(B_F)$ sur $H^{q}(X)$.
\item[--] Pour un groupe $Y$ et un sous-groupe $F$, la fibration
$Y/F \to B_F \to B_Y$ : $H^{p}(B_Y, \{H^{q}(Y/F)\}) \Rightarrow
H^{p+q}(B_F)$, le système local étant défini par l'action de $\pi_0(Y)$ sur
$H^{q}(Y/F)$.
\item[--] \emph{Hochschild--Serre.} Pour un groupe $X$, un sous-groupe
distingué $F$ et $Y = X/F$, la fibration $B_F \to B_X \to B_Y$ :
$H^{p}(B_Y, H^{q}(B_F)) \Rightarrow H^{p+q}(B_X)$, avec coefficients locaux ;
pour des groupes discrets, c'est
$H^{p}(Y, H^{q}(F)) \Rightarrow H^{p+q}(X)$.
\end{enumerate}
La page met en marge, en face de la dernière, « Leray-Čech » à côté de
« Hochschild-Serre » : la suite (B) de la page 19, appliquée à
$G/F' \to G/F$ avec $F'$ distingué, est son analogue discret.\footnote{Pour
Leray, la page écrit $H^{p}(Y, H^{q}(F))$ sans système local : c'est exact si
$Y$ est simplement connexe, ou si $\pi_1(Y)$ opère trivialement sur la
cohomologie des fibres ; en général il faut les coefficients locaux, qu'elle
met d'ailleurs dans les deux suites suivantes. Dans la troisième, elle écrit
$X$ là où elle vient de nommer le groupe $Y$ (le $B_X$ est biffé, le
$\pi_0(X)$ du N.B. ne l'est pas) ; dans la quatrième, l'aboutissement est
écrit $H^{*}(B_G)$ pour $H^{*}(B_X)$. On corrige. Trois lignes serrées entre
la troisième et la quatrième, sur le cas où $Y$ est discret, sont
illisibles sauf leurs symboles.}

\pagerange{25}{27}
\section*{IV. Familles limitées (pages 25 à 43)}

La page 25 est une couverture : « Familles limitées », et
« cf. SGA 6 XIII ».\footnote{L'exposé XIII de SGA 6 (1966-1967), de
S.~Kleiman, « Les théorèmes de finitude pour le foncteur de Picard », traite
des familles limitées de faisceaux ; la couverture y renvoie, sans qu'on
puisse dire si ces feuillets le préparent ou le suivent.}

\emph{Définition.} Soit $X$ de type fini sur $S$ noethérien. Une famille
$(F_i)$ de faisceaux cohérents, $F_i$ vivant sur la fibre géométrique
$X_{s_i} \otimes_{k(s_i)} K_i$, est \emph{limitée} (aujourd'hui :
\emph{bornée}) s'il existe un $S$-schéma $S'$ de type fini et un faisceau
cohérent $F$ sur $X \times_S S'$ tels que chaque $F_i$ soit, à extension du
corps près, isomorphe à une fibre $F_{s'}$ de $F$.\footnote{La définition
n'est pas écrite dans le dossier ; c'est celle que supposent les pages 34 et
43 (« familles universelles »), et celle de SGA 6, XIII. Le sens précis de
« à extension du corps près » est l'objet des pages 28 à 33.}

\subsection*{Les quotients d'un faisceau fixe (pages 26 et 27)}

L'énoncé visé est le suivant : \emph{sur $\mathbf{P}^{n}_{S}$, les
faisceaux $F_i$ quotients d'un faisceau cohérent fixé $E$, et dont les
polynômes de Hilbert $P_{F_i}$ sont en nombre fini, forment une famille
limitée.} L'énoncé est vrai ; c'est la bornitude qui sous-tend
l'existence du schéma $\mathrm{Quot}$.\footnote{Il découle de la
projectivité du schéma $\mathrm{Quot}$, construit par Grothendieck dans les
\emph{Techniques de construction} (FGA, exposé 221, 1961), ou de la
régularité de Castelnuovo--Mumford (Mumford, 1966), qui est la voie de
SGA 6, XIII. On ne prétend pas que les pages 26 et 27 en contiennent une
preuve.}

La preuve esquissée procède par récurrence sur
$d = \sup \dim \operatorname{Supp} F_i$, le cas $d < 0$ étant trivial. Pour
$d \geq 0$, on écrit
\[
  0 \to F''_i \to F_i \to F'_i \to 0,
\]
où $\dim \operatorname{Supp} F''_i < d$ et où $F'_i$ est contrôlé par les
cycles de dimension $d$ associés aux $F_i$, de degré borné. Ces cycles
forment une famille limitée par un « résultat de Chow » : \emph{les cycles
de degré borné de $\mathbf{P}^{n}_{\mathbf{Z}}$ forment une famille
limitée}, c'est-à-dire sont les fibres d'un
$\Gamma \subset \mathbf{P}^{n}_{T}$ au-dessus d'un $T$ de type fini sur
$\mathbf{Z}$. Les $F'_i$ sont alors limités, leurs polynômes de Hilbert sont
en nombre fini, donc aussi ceux des $F''_i$, puisque
$P_{F''_i} = P_{F_i} - P_{F'_i}$ ; et l'on conclut par récurrence que les
$F''_i$, puis les $F_i$, sont limités.\footnote{Pages 26 et 27, d'écriture
très rapide : seuls les signes et quelques mots se lisent. Les étapes
cerclées (d), (e), (f) de la marge sont communes aux deux pages, qui
reprennent le même plan. La page 27 dit des $F''_i$ qu'ils sont
« quotients d'un » faisceau dont le nom est illisible, puis « j'y
reviendrai » : l'étape qui ramène les $F''_i$ à l'hypothèse de récurrence
(ce sont des sous-faisceaux des $F_i$, non des quotients de $E$) n'est pas
lisible, et on ne la reconstruit pas. Le « résultat de Chow » est ce qu'on
appelle aujourd'hui l'existence (et la quasi-projectivité) de la variété de
Chow ; la page l'énonce sur $\mathbf{Z}$.}

\pagerange{28}{29}
\subsection*{Faisceaux sur des extensions du corps de base (pages 28 et 29)}

\emph{Définition 1.} Soient $X$ un schéma sur un corps $k$, $K_1$, $K_2$ deux
extensions de $k$, et $F_i$ un faisceau cohérent sur $X \otimes_k K_i$. On
dit que $F_1$ et $F_2$ sont \emph{$k$-équivalents} s'il existe une
extension $K'$ de $k$ et des $k$-plongements $u_i : K_i \to K'$ tels que les
faisceaux $F_i \otimes_{K_i} K'$ sur $X \otimes_k K'$ soient isomorphes.

C'est une relation d'équivalence sur les couples $(K, F)$ : pour la
transitivité, deux extensions d'un même corps se plongent toujours dans une
extension commune.\footnote{La justification de la transitivité est en
partie illisible sur la page ; on la complète.}

\emph{Proposition 1} ($X$ de type fini sur $k$). Tout faisceau cohérent sur
$X \otimes_k K'$ est équivalent à un faisceau défini sur
$X \otimes_k K$, $K \subset K'$ de type fini sur $k$. Si $F_1$ et $F_2$ sont
équivalents et définis sur des extensions de type fini, on peut, dans la
définition 1, prendre $K'$ de type fini.

Par suite, les classes d'équivalence forment un ensemble, en bijection avec
les classes de faisceaux définis sur des extensions de type fini de $k$ ; on
le note $\mathfrak{F}(X/k)$. Il dépend de $k$ de façon
contravariante.\footnote{Pages 28 et 29. Le premier corollaire est dit
résulter de « Chap. V, § 1 », renvoi qu'on n'identifie pas. La page pose
ensuite une « Question » : peut-on prendre pour $K'$ une extension
universelle ? La page 30 y répond, dans la lecture qu'on en donne. Un
bloc barré de trois diagonales, sur un changement de corps de base, est
omis.}

\pagerange{30}{30}
\subsection*{Une extension universelle (page 30)}

\emph{Corollaire 5.} Soient $\widetilde{K}$ une extension algébriquement
close de $k$ de degré de transcendance infini, et
$\widetilde{X} = X \otimes_k \widetilde{K}$. Alors deux faisceaux cohérents
$F$, $G$ sur $\widetilde{X}$ sont $k$-équivalents si et seulement s'il
existe $\sigma \in \operatorname{Aut}(\widetilde{K}/k)$ tel que
$F \simeq \sigma_{\widetilde{X}}(G)$, où $\sigma_{\widetilde{X}}$ est
l'automorphisme $\mathrm{id}_X \times \sigma$ de $\widetilde{X}$ ; et
$\mathfrak{F}(X/k)$ s'identifie au quotient de l'ensemble des classes
d'isomorphisme de faisceaux cohérents sur $\widetilde{X}$ par
$\operatorname{Aut}(\widetilde{K}/k)$.

Les deux hypothèses sur $\widetilde{K}$ servent ainsi : toute extension de
type fini de $k$ s'y plonge, ce qui donne la surjectivité ; et deux
plongements d'un même sous-corps de type fini diffèrent d'un automorphisme de
$\widetilde{K}$ (on prolonge l'isomorphisme à des bases de transcendance,
de même cardinal infini, puis aux clôtures algébriques), ce qui
donne l'injectivité.\footnote{La page écrit « une extension
algébrique[ment close] de $k$ … transcendance … », le mot « close » et la
condition sur le degré de transcendance étant d'une lecture incertaine ou
illisible. On complète par les deux hypothèses qui rendent l'énoncé vrai, et
l'argument est le nôtre. Le numéro « 5 » du corollaire est incertain.}

\pagerange{31}{33}
\subsection*{Opérations sur les classes (pages 30 à 33)}

Supposons donnés $X$, $Y$ de type fini sur $k$ et, pour toute extension
algébriquement close $K$ de $k$ et tous faisceaux cohérents
$F_1, \ldots, F_n$ sur $X_K$, un ensemble $T(F_1, \ldots, F_n)$ de classes de
faisceaux cohérents sur $Y_K$, de sorte que :
\begin{enumerate}
\item[(i)] $T(F_1, \ldots, F_n)$ ne change pas si l'on remplace les $F_i$
par des faisceaux isomorphes ;
\item[(ii)] $T$ est compatible à l'extension du corps de base : pour
$K \subset K'$, $T(F_1 \otimes_K K', \ldots, F_n \otimes_K K')$ est, à
isomorphisme près, l'ensemble des $G \otimes_K K'$, $G \in T(F_1, \ldots,
F_n)$.
\end{enumerate}
Pour des parties $A_1, \ldots, A_n$ de $\mathfrak{F}(X/k)$, on note
$\mathfrak{F}(A_1, \ldots, A_n) \subset \mathfrak{F}(Y/k)$ l'ensemble des
classes des éléments des $T(F_1, \ldots, F_n)$, où les $F_i$ sont des
faisceaux sur $\widetilde{X}$ de classes dans les $A_i$ ; la définition ne
dépend pas du choix de $\widetilde{K}$.

Si l'on part de faisceaux $F_i$ définis sur des extensions \emph{distinctes}
$K_i$, de classes $\xi_i$, on obtient $\mathfrak{F}(\xi_1, \ldots, \xi_n)$
en parcourant \emph{toutes} les extensions composées $L_\alpha$ des $K_i$
sur $k$ et en réunissant les classes des
$T(F_1 \otimes L_\alpha, \ldots, F_n \otimes L_\alpha)$. Les classes $\xi_i$
ne fixent pas la position relative des corps $K_i$, et chaque manière de
les composer donne en général d'autres faisceaux.\footnote{Pages 30 à 32,
d'écriture rapide et en grande partie illisible ; on en garde les énoncés.
L'ensemble $T$ est introduit au bas de la page 30 et la page 31 le
prolonge ; « Déf. », encadré, y est biffé. La page 32 écrit
$\mathfrak{F}(F_1, \ldots, F_n)$ pour $T(F_1, \ldots, F_n)$. Les « types
d'extensions composées » reviennent page 34, où ils sont les points d'un
produit fibré.}

\emph{Exemples} (page 33) : a) les foncteurs $\mathcal{H}om$, $\otimes$ et
leurs dérivés $\mathcal{T}or$, $\mathcal{E}xt$ (deux arguments, $X = Y$) ;
(d) un ensemble $N(F, G)$ attaché à deux faisceaux ; (b) et (c) deux
opérations à un argument attachées à un morphisme $f : X \to Y$ donné, la
seconde « et ses foncteurs dérivés » ; (e) l'ensemble des extensions de $F$
par $G$, c'est-à-dire des faisceaux $E$ qui figurent au milieu d'une suite
exacte $0 \to G \to E \to F \to 0$ ; (f) une opération attachée à une
filtration décroissante de $F$ ; (g) le faisceau
$\mathcal{O}_{(\operatorname{Supp} F)_{\mathrm{red}}}$ ; (h) les composantes
primaires de $F$ associées aux composantes irréductibles de son support.
Pour les deux dernières il faut $K$ algébriquement clos : être réduit, être
irréductible ne sont pas stables par extension inséparable ou non
algébrique du corps, et (ii) serait faux.\footnote{La page numérote, dans
cet ordre, a, (d), (b), (c), (e), (f), (f), (h) ; on a rebaptisé (g) le
second (f). Les noms des opérations (b), (c), (d) et du premier (f) sont
illisibles ; on ne les devine pas, sinon pour noter que le plan de la page
43 cite images inverses et images directes avec leurs dérivés $R^{p}f_{*}$,
ce qui conviendrait à (b) et (c) --- rapprochement qui est le nôtre. La
justification par la non-stabilité de la réduction et de l'irréductibilité
est aussi la nôtre ; la page se contente d'écrire « $K$ alg. clos ». La
ligne (i), biffée, est omise.}

\pagerange{34}{35}
\subsection*{Stabilité des familles limitées (pages 34 à 37)}

On note maintenant $\mathfrak{F}(X/S)$ la réunion des $\mathfrak{F}(X_s/k(s))$,
$s \in S$, et, pour un faisceau cohérent $F$ sur $X \times_S S'$, $S'$ de type
fini sur $S$, $\mathcal{E}_{X/S}(F) \subset \mathfrak{F}(X/S)$ l'ensemble
des classes de ses fibres $F_{s'}$. Une partie de $\mathfrak{F}(X/S)$ est
limitée si elle est contenue dans un $\mathcal{E}_{X/S}(F)$.

\emph{(a) Si $A, B \subset \mathfrak{F}(X/S)$ sont limités, $A \otimes B$,
$\mathcal{H}om(A, B)$, $\mathcal{T}or_i(A, B)$ et $\mathcal{E}xt^{i}(A, B)$
le sont.}

\emph{Le produit tensoriel.} Soient $A \subset \mathcal{E}(F)$, $F$ sur
$X \times_S S'$, et $B \subset \mathcal{E}(G)$, $G$ sur $X \times_S S''$.
Posons $S''' = S' \times_S S''$ et $H = F \otimes G$ (images inverses) sur
$X \times_S S'''$. Tout type d'extension composée de $k(s')$ et $k(s'')$ sur
$k(s)$ provient d'un point $s''' \in S'''$ au-dessus de $s'$ et $s''$ --- ce
sont les points du produit fibré (EGA I, § 3.4) ---, et
\[
  (F_{s'} \otimes_{k(s')} K) \otimes_{X_{s} \otimes_{k(s)} K}
  (G_{s''} \otimes_{k(s'')} K) \simeq H_{s'''} \otimes_{k(s''')} K .
\]
Comme $S'''$ est de type fini sur $S$, on a
$\mathcal{E}_{X/S}(F) \otimes \mathcal{E}_{X/S}(G) \subset
\mathcal{E}_{X/S}(H)$. Le produit tensoriel commute à tout changement de
base, et rien d'autre n'est nécessaire.\footnote{La page écrit $H_{s'''}$
sans extension du corps dans le second membre ; on l'ajoute. Elle dessine
en marge le cube des $X$, $X'$, $X''$, $X'''$ au-dessus des $S$, $S'$,
$S''$, $S'''$.}

\emph{Les $\mathcal{H}om$.} Ceux-ci ne commutent pas au changement de base,
et il faut stratifier. Soient $\bar{F}$, $\bar{G}$ les images inverses de $F$
et $G$ sur $X \times_S S'''$. Il existe une partition finie de $S'''$ en
parties localement fermées $T_i$, munies de leur structure réduite, telle
que, en posant $F_i = \bar{F} \otimes T_i$, $G_i = \bar{G} \otimes T_i$, la
formation de $\mathcal{H}om_{X_{T_i}}(F_i, G_i)$ commute à tout changement
de base $T' \to T_i$. En effet, localement $\bar{F}$ a une présentation
$L_1 \to L_0 \to \bar{F} \to 0$ par des modules libres de type fini, et
\[
  \mathcal{H}om(\bar{F}, \bar{G}) =
  \operatorname{Ker}\bigl(\mathcal{H}om(L_0, \bar{G}) \to
  \mathcal{H}om(L_1, \bar{G})\bigr) ;
\]
par platitude générique, sur un ouvert dense de $(S''')_{\mathrm{red}}$,
$\bar{G}$ et le conoyau de cette flèche sont plats, et le noyau d'une flèche
entre modules plats, de conoyau plat, commute à tout changement de base. La
récurrence noethérienne fournit la partition. Les $\mathcal{T}or_i$ et les
$\mathcal{E}xt^{i}$ se traitent de même.\footnote{Page 34, la page écrit
$\bar{G} = G \otimes_{S'} S'''$, où il faut $\otimes_{S''}$, comme dans
$G \otimes_{S''} T_i$ qui suit. Page 35, le noyau est écrit
$\operatorname{Ker}(\mathcal{H}om(L_0, F) \to \mathcal{H}om(L_1, F))$ : on
lit $\bar{G}$ pour $F$, puisque c'est $\mathcal{H}om(\bar{F}, \bar{G})$
qu'on calcule. L'argument de platitude, entre crochets sur la page, est
presque entièrement illisible ; celui qu'on donne est le nôtre, et il est
standard (platitude générique, EGA IV, 6.9.1). La phrase sur les
$\mathcal{T}or_i$ et $\mathcal{E}xt^{i}$ est illisible sauf ces mots.}

\pagerange{35}{36}
\emph{Lemme} (pages 35 et 36). Soient $T$ noethérien, $Y$ de type fini sur
$T$, $K^{\bullet}$ un complexe de faisceaux cohérents sur $Y$, de
différentielle de degré $+1$, et $N$ un entier. Il existe une partition
finie de $T$ en parties localement fermées $T_i$, munies de leur structure
réduite, telle qu'en posant $K^{(i)} = K \times_T T_i$, les homomorphismes
canoniques
\[
  H^{q}\bigl(K^{(i)}\bigr) \otimes_{T_i} T' \longrightarrow
  H^{q}\bigl(K^{(i)} \otimes_{T_i} T'\bigr)
\]
soient des isomorphismes pour $|q| \leq N$ et tout $T' \to T_i$.

\emph{Preuve.} On peut remplacer $T$ par $T_{\mathrm{red}}$. Par platitude
générique, il existe un ouvert dense $U$ de $T$ sur lequel les $K^{i}$, les
cobords $B^{i}$, les cocycles $Z^{i}$ et les $H^{i}$ sont plats pour
$|i| \leq N + 1$ ; alors les suites exactes $0 \to Z^{i} \to K^{i} \to
B^{i+1} \to 0$ et $0 \to B^{i} \to Z^{i} \to H^{i} \to 0$ restent exactes
après tout changement de base, et la cohomologie commute au changement de
base en ces degrés. On conclut par récurrence noethérienne sur
$T \setminus U$.\footnote{La preuve de la page est presque entièrement
illisible : on y lit « nilpotent », « réduit », un ouvert $U$, et la suite
$K^{i-1} \to K^{i} \to K^{i+2}$, où l'on lit $K^{i+1}$. Celle qu'on donne
est la nôtre. La note marginale renvoie à « IV, § 3 », renvoi qu'on
n'identifie pas.}

\pagerange{36}{37}
\emph{(b) Images inverses.} Soit $f : X \to Y$ un $S$-morphisme entre
$S$-schémas de type fini. Si $A \subset \mathfrak{F}(Y/S)$ est limité,
$f^{*}(A)$ et les $\mathcal{T}or^{\mathcal{O}_Y}_{*}(A, \mathcal{O}_X)$ le
sont. De même, pour $X$ et $Z$ au-dessus de $Y$,
$A \subset \mathfrak{F}(X/S)$ et $B \subset \mathfrak{F}(Z/S)$ limités, les
$\mathcal{T}or^{\mathcal{O}_Y}_{i}(A, B)$, faisceaux sur $X \times_Y Z$, le
sont.\footnote{Page 36, l'indice de ce dernier $\mathcal{T}or$ se lit
$\mathcal{O}_Z$ ; c'est $\mathcal{O}_Y$ qu'il faut. L'énoncé est
introduit par « Il y a un énoncé analogue », la suite étant illisible ; on
l'écrit dans la forme qu'appelle le contexte.}

La page 37 pose alors la question qui commande le reste du dossier. Pour
des $X_i$ au-dessus de $Y$ ($1 \leq i \leq n$), des $F_i$ sur les $X_i$, et
un changement de base $S' \to S$, il y a un homomorphisme canonique
\[
  \mathcal{T}or^{\mathcal{O}_Y}_{*}(F_1, \ldots, F_n) \otimes_S S'
  \longrightarrow
  \mathcal{T}or^{\mathcal{O}_{Y'}}_{*}(F'_1, \ldots, F'_n),
\]
où $\mathcal{T}or^{\mathcal{O}_Y}_{*}(F_1, \ldots, F_n)$ est le
$\mathcal{T}or$ multiple, faisceau sur $X_1 \times_Y \cdots \times_Y X_n$.
« À quelle condition est-ce un isomorphisme ? » Il suffit que les
$\mathcal{T}or^{\mathcal{O}_Y}_{q}(F_1, \ldots, F_n)$ soient plats sur
$S$ --- c'est l'objet des pages qui suivent.\footnote{La page écrit
« $Y$-plats ». Le changement de base est pris sur $S$ (la marge aligne
$X_1\ X_2$, $X'_1\ X'_2$, $Y$, $Y'$, $S$, $S'$), et c'est la platitude sur
$S$ que le lemme de la page 38 demande ; on lit donc « $S$-plats ». Une
première formule, avec un $\mathrm{Tor}^{\mathcal{O}_S}_{p}$ de
$\mathrm{Tor}^{\mathcal{O}_Y}_{q}$, est barrée.}

\pagerange{38}{40}
\subsection*{Changement de base pour les $\mathcal{T}or$ (pages 38 à 40)}

\emph{Proposition.} Soient $S$ noethérien, $Y$ de type fini sur $S$, des
$X_i$ ($1 \leq i \leq n$) de type fini sur $Y$, des $F_i$ cohérents sur les
$X_i$, et $N$ un entier. Il existe un ouvert dense $U$ de $S$ tel qu'en
posant $S' = U_{\mathrm{red}}$, $X'_i = X_i \times_S S'$,
$Y' = Y \times_S S'$, $F'_i = F_i \otimes_S S'$, les homomorphismes
canoniques
\[
  \mathcal{T}or^{\mathcal{O}_{Y'}}_{p}(F'_1, \ldots, F'_n) \otimes_{S'} T
  \longrightarrow
  \mathcal{T}or^{\mathcal{O}_{Y_T}}_{p}(F'_{1T}, \ldots, F'_{nT})
\]
soient des isomorphismes pour $p \leq N$ et tout changement de base
$T \to S'$.

Elle résulte, par platitude générique (qui rend plats sur $U$ le schéma $Y$,
les $X_i$, les $F_i$ et les $\mathcal{T}or_q$ pour $q \leq N$), du

\emph{Lemme.} Supposons $Y$ et les $X_i$ plats sur $S$, les $F_i$ plats sur
$S$, et les $\mathcal{T}or^{\mathcal{O}_Y}_{q}(F_1, \ldots, F_n)$ plats sur
$S$ pour $q \leq N$. Alors pour tout $T \to S$ l'homomorphisme canonique
\[
  \mathcal{T}or^{\mathcal{O}_Y}_{q}(F_1, \ldots, F_n) \otimes_S T
  \longrightarrow
  \mathcal{T}or^{\mathcal{O}_{Y_T}}_{q}(F_{1(T)}, \ldots, F_{n(T)})
\]
est un isomorphisme pour $q \leq N$.\footnote{Page 38. La Proposition
écrit les $\mathrm{Tor}$ sur $\mathcal{O}_{S'}$ et $\mathcal{O}_T$, et le
Lemme sur $\mathcal{O}_S$ dans la flèche (sur $\mathcal{O}_Y$ dans les
hypothèses) : ce sont des $\mathcal{T}or$ relatifs à $Y$, comme le montre la
preuve, et on les écrit ainsi. Le Lemme de la page ne demande pas les $F_i$
plats sur $S$ ; sa preuve (page 39) s'en sert --- pour que
$L^{(i)} \otimes_A A'$ résolve $M_i \otimes_A A'$ ---, le mot « plat » y
étant d'ailleurs écrit puis biffé après $M_i$, et le lemme analogue de la
page 41 pour les $\mathcal{E}xt$ le demande. On l'ajoute, sans affirmer que
l'énoncé sans elle est faux. Dans la Proposition, elle est assurée par la
platitude générique. La marge porte « Faut-il … ? » en face du Lemme, et le
mot « Important » entre parenthèses est d'une lecture incertaine.}

\emph{Preuve} (pages 39 et 40). On se ramène à $S$, $Y$, les $X_i$, $T$
affines, d'anneaux $A$, $B$, $B_i$, $A'$, les $F_i$ étant définis par des
$B_i$-modules $M_i$. Soit $L^{(i)}$ une résolution libre de $M_i$ par des
$B_i$-modules, et
\[
  K = \bigotimes_{B,\, 1 \leq i \leq n} L^{(i)}
\]
(produit tensoriel sur $B$ des complexes), de sorte que
$H_{q}(K) = \mathrm{Tor}^{B}_{q}(M_1, \ldots, M_n)$. Les $L^{(i)}$ sont
plats sur $A$, puisque les $B_i$ le sont ; comme $M_i$ est plat sur $A$,
$L^{(i)\prime} = L^{(i)} \otimes_A A'$ est une résolution de
$M'_i = M_i \otimes_A A'$ comme $B'_i$-module, $B'_i = B_i \otimes_A A'$.
D'où
\[
  K' = K \otimes_A A' = \bigotimes_{B',\, 1 \leq i \leq n} L^{(i)\prime},
  \qquad H_{*}(K') = \mathrm{Tor}^{B'}_{*}(M'_1, \ldots, M'_n).
\]
Comme $K$ est un complexe (borné inférieurement) de $A$-modules plats, on a
la suite spectrale de Künneth
\[
  E^{2}_{p,q} = \mathrm{Tor}^{A}_{p}\bigl(H_{q}(K), A'\bigr) \Longrightarrow
  H_{p+q}(K \otimes_A A'),
\]
c'est-à-dire
\[
  \mathrm{Tor}^{A}_{p}\bigl(\mathrm{Tor}^{B}_{q}(M_1, \ldots, M_n), A'\bigr)
  \Longrightarrow \mathrm{Tor}^{B'}_{p+q}(M'_1, \ldots, M'_n).
\]
Si les $\mathrm{Tor}^{B}_{q}(M_1, \ldots, M_n)$ sont plats sur $A$ pour
$q \leq N$, les termes $E^{2}_{p,q}$ avec $p > 0$ et $q \leq N$ sont nuls ;
en degré total $n \leq N$ il ne reste que $E^{2}_{0,n}$, qu'aucune
différentielle ne touche, et
\[
  \mathrm{Tor}^{B}_{p}(M_1, \ldots, M_n) \otimes_A A' \longrightarrow
  \mathrm{Tor}^{B'}_{p}(M'_1, \ldots, M'_n)
\]
est un isomorphisme pour $p \leq N$.\footnote{La page écrit la suite
spectrale sous la forme $H_{*}(K \otimes_A A') \Leftarrow
\mathrm{Tor}^{A}_{p}(H_q(K), A')$, puis l'applique ; l'analyse des
différentielles en degré $\leq N$ est la nôtre. Dans la chaîne
$K' = \ldots$, un terme intermédiaire est noirci ; l'indice du dernier
$\mathrm{Tor}^{\cdot}_{q}$ de la page 39 est illisible, on lit $B$.}

\emph{Corollaire} (page 40). Il existe une partition finie de $S$ en
parties localement fermées $S_\alpha$ (structure réduite) telle qu'en posant
$X^{(\alpha)}_i = X_i \times_S S_\alpha$, $Y^{(\alpha)} = Y \times_S S_\alpha$,
$F^{(\alpha)}_i = F_i \otimes_S S_\alpha$, l'homomorphisme canonique
\[
  \mathcal{T}or^{\mathcal{O}_{Y^{(\alpha)}}}_{p}\bigl(F^{(\alpha)}_1, \ldots,
  F^{(\alpha)}_n\bigr) \otimes_{S_\alpha} T \longrightarrow
  \mathcal{T}or^{\mathcal{O}_{Y_T}}_{p}\bigl(F^{(\alpha)}_{1(T)}, \ldots,
  F^{(\alpha)}_{n(T)}\bigr)
\]
soit un isomorphisme pour $p \leq N$ et tout $T \to S_\alpha$. C'est la
Proposition et la récurrence noethérienne.\footnote{La page écrit les
$\mathrm{Tor}$ sur $\mathcal{O}_{S_\alpha}$ et « $|p| \leq N$ » ; les
$\mathcal{T}or$ n'ont que des indices $p \geq 0$.}

\pagerange{40}{42}
\subsection*{Changement de base pour les $\mathcal{E}xt$ (pages 40 à 42)}

\emph{Proposition} (page 40). Soient $X$ de type fini sur $S$ noethérien,
$F$ et $G$ cohérents sur $X$, $N$ un entier. Il existe un ouvert dense $U$
de $S$ tel qu'en posant $S' = U_{\mathrm{red}}$, $X' = X \times_S S'$,
$F' = F \otimes_S S'$, $G' = G \otimes_S S'$, l'homomorphisme canonique
\[
  \mathcal{E}xt^{p}_{\mathcal{O}_{X'}}(F', G') \otimes_{S'} T \longrightarrow
  \mathcal{E}xt^{p}_{\mathcal{O}_{X_T}}\bigl(F'_{(T)}, G'_{(T)}\bigr)
\]
soit un isomorphisme pour $p \leq N$ et tout $T \to S'$.\footnote{La page
écrit $\mathcal{E}xt^{p}_{\mathcal{O}_{S'}}$ et
$\mathcal{E}xt^{p}_{\mathcal{O}_{T}}$, où il faut les faisceaux d'anneaux de
$X'$ et $X_T$ ; « pour tout changement de base $T \to S$ », où il faut
$T \to S'$ ; « $|p| \leq N$ » ; et $F \times_S S'$ pour
$F \otimes_S S'$.}

Elle résulte par platitude générique du lemme suivant, écrit « avec les
notations précédentes, mais sans conditions de finitude ».

\emph{Lemme} (pages 41 et 42). Supposons $X$ plat sur $S$, $F$ et $G$ plats
sur $S$, et les $\mathcal{E}xt^{i}(F, G)$ plats sur $S$ pour $i \leq N$.
Supposons de plus que tout point de $X$ ait un voisinage $U$ sur lequel
$F|_U$ possède une résolution $L_{*} \to F|_U$ de longueur $\leq N + 1$ par
des modules localement libres de type fini, telle que le conoyau de
$\mathcal{H}om(L_N, G) \to \mathcal{H}om(L_{N+1}, G)$ soit plat sur $S$.
Alors pour tout $T \to S$ les homomorphismes canoniques
\[
  \mathcal{E}xt^{i}_{\mathcal{O}_X}(F, G) \otimes_S T \longrightarrow
  \mathcal{E}xt^{i}_{\mathcal{O}_{X_T}}\bigl(F_{(T)}, G_{(T)}\bigr)
\]
sont des isomorphismes pour $i \leq N$.\footnote{Page 41. La seconde
hypothèse est une longue addition écrite en oblique dans la marge gauche,
sans signe de renvoi sinon le trait qui borde l'énoncé ; son bloc central est
biffé, et on n'en garde que ce qui reste. Il faut entendre « résolution
exacte en degrés $\leq N$ », comme le précise la preuve. La page écrit
$\mathcal{E}xt^{i}_{\mathcal{O}_X}(F, G) \times_S T$.}

\emph{Preuve.} La définition des homomorphismes est locale, et l'on se
ramène à $S$, $X$, $T$ affines, d'anneaux $A$, $B$, $A'$ ; on laisse au
lecteur de vérifier que les isomorphismes obtenus se recollent. Soient $M$,
$N$ les $B$-modules qui définissent $F$, $G$ ; $B$, $M$ et $N$ sont plats
sur $A$. Posons $B' = B \otimes_A A'$, $M' = M \otimes_A A'$,
$N' = N \otimes_A A'$. Soit $L_{*}$ un complexe de $B$-modules libres de type
fini en degrés $\leq N + 1$, acyclique en degrés $\leq N$, augmenté vers
$M$. Comme $M$ et les $L_i$ sont plats sur $A$, $L_{*} \otimes_A A'
= L_{*} \otimes_B B'$ est encore, en degrés $\leq N$, une résolution de
$M'$, formée de libres de type fini. Donc, pour $p \leq N$,
\[
  \operatorname{Ext}^{p}_{B'}(M', N') = H^{p}\bigl(\operatorname{Hom}_{B'}(L'_{*},
  N')\bigr) = H^{p}\bigl(\operatorname{Hom}_{B}(L_{*}, N) \otimes_A A'\bigr).
\]
Soit $K^{*} = \operatorname{Hom}_{B}(L_{*}, N)$. On a donc, pour $p \leq N$,
\[
  \operatorname{Ext}^{p}_{B'}(M', N') = H^{p}(K^{*} \otimes_A A'), \qquad
  \operatorname{Ext}^{p}_{B}(M, N) = H^{p}(K^{*}).
\]
Les $K^{i}$ sont plats sur $A$ (car $N$ l'est et les $L_i$ sont libres de
type fini), les $H^{p}(K^{*})$ aussi pour $p \leq N$ par hypothèse, et le
conoyau $C^{N+1}$ de $K^{N} \to K^{N+1}$ est plat sur $A$. Il s'ensuit, en
descendant à partir de $C^{N+1}$ le long des suites exactes
$0 \to H^{p} \to C^{p} \to K^{p+1} \to C^{p+1} \to 0$, que les cobords, les
cocycles et les conoyaux $C^{p}$ sont plats pour $p \leq N$, donc que
\[
  H^{p}(K^{*}) \otimes_A A' \longrightarrow H^{p}(K^{*} \otimes_A A')
\]
est un isomorphisme pour $p \leq N$ ; c'est-à-dire
\[
  \operatorname{Ext}^{p}_{B'}(M', N') \simeq \operatorname{Ext}^{p}_{B}(M, N)
  \otimes_A A' \qquad (p \leq N),
\]
ce qu'il fallait démontrer.\footnote{Page 42. L'indice de $\operatorname{Hom}$, deux fois, se lit $A$
là où il faut $B$. La page dit les $K^{i}$ « $A$-libres » (lecture
incertaine) : ils sont $A$-plats, ce qui suffit. Elle conclut que
l'homomorphisme de changement de base est un isomorphisme « en tt
dimension », puis « en particulier pour $p \leq N$ » : les hypothèses ne
donnent que $p \leq N$, et c'est ce qu'on énonce. La descente par les suites
exactes est notre détail ; la page dit « Il s'ensuit ». Le $N$ qui nomme le
module et le $N$ qui borne les degrés sont deux lettres identiques sur la
page ; on les garde, le contexte les distingue.}

\pagerange{43}{43}
\subsection*{Le plan d'un chapitre (page 43)}

La dernière page, d'une écriture posée, donne le plan d'un chapitre sur les
familles limitées de faisceaux cohérents et de complexes de faisceaux
cohérents, par les seuls intitulés de ses énoncés :
\begin{enumerate}
\item[Prop. 1] une réunion finie de familles limitées est limitée ;
\item[Prop. 2] dans la définition, on peut supposer les « familles
universelles » plates sur la base --- ce que donnent la platitude générique
et une stratification, comme pages 34 à 40 ;
\item[Prop. 3] un changement de base $S' \to S$ transforme famille limitée
en famille limitée, avec une réciproque si l'image de $S' \to S$ contient
tous les points de $S$ où vivent les membres de la famille ;
\item[Prop. 4] les faisceaux de cohomologie $\mathcal{H}^{p}(K)$ d'une
famille limitée de complexes ;
\item[Prop. 5] une cohomologie relative à plusieurs arguments,
$H^{p}_{f_1, \ldots, f_n}(K_1, \ldots, K_n)$, pour des $f_i : X_i \to Z_i$
propres ;
\item[Cor. 1--3] les $\mathcal{T}or$ locaux, les images inverses par
$X' \to X$, les images directes et les $R^{p}f_{*}$ ;
\item[Prop. 6] les $\mathcal{E}xt^{p}_{\mathcal{O}_X}(K, L)$, et en
corollaire les $\mathcal{H}om(F, G)$ ;
\item[Prop. 7] des $\mathcal{E}xt^{p}$ relatifs à $f$, « quand ça marche »,
et en corollaire leur $\mathcal{E}xt^{1}$ « en tout cas ».
\end{enumerate}
Aucun de ces énoncés n'est démontré dans le dossier ; les pages 34 à 42 en
fournissent les outils pour les Propositions 2, 4 et 6.\footnote{L'indice de
Prop. 5 est lu $H^{s}_{f_1, \ldots, f_n}{}^{p}$ ; celui des
$\mathcal{E}xt$ de Prop. 7 se lit « $X$, $f$ », d'une lecture incertaine.
Ce sont, dans le vocabulaire actuel, des $\mathcal{E}xt$ relatifs au sens
d'EGA III, 7.7, ou d'Altman--Kleiman ; on ne prétend pas que ce soit ce que
la page désigne. Le crochet ouvert après Prop. 6 reste sans suite, et
l'indice $\mathcal{O}_X$ peut appartenir au corollaire.}

\end{document}
